\documentclass[10pt,a4paper]{article}
\usepackage{../../styles/neet_style}

\pagestyle{fancy}
\fancyhf{}
\fancyhead[L]{\small \textbf{NEET UG 2027: Optics \& Modern Physics Master Theory}}
\fancyhead[R]{\small Satya Prakash}
\fancyfoot[L]{\small \textbf{NEET UG 2027 Prep Series}}
\fancyfoot[C]{\thepage}
\fancyfoot[R]{\small \textbf{Physics Chapter Notes}}

\begin{document}

% Title Header Banner
\begin{center}
  {\Huge \textbf{NEET UG 2027: Master Chapter Theory}}\\[6pt]
  {\Large \textbf{Ray Optics, Wave Optics, Quantum Radiation, Atomic \& Nuclear Physics, Semiconductors}}\\[4pt]
  {\small \textbf{Prepared, Compiled \& Authored by Satya Prakash} $\cdot$ Strictly Calibrated to Revised NMC Syllabus}
  \vspace{4pt}
  \hrule height 1.5pt
\end{center}

\section{Ray Optics, Reflection, Refraction \& Total Internal Reflection}

\begin{conceptbox}[Foundations of Geometric Optics \& Sign Conventions]
\textbf{1. Cartesian Sign Convention (Strict Standard):}
\begin{itemize}
  \item All optical distances are measured along the principal axis starting from the optical pole ($P$) or optical centre ($O$).
  \item Distances measured in the direction of incident light are \textbf{positive} ($+$); distances opposite to incident light are \textbf{negative} ($-$).
  \item Heights measured perpendicularly upwards above the principal axis are \textbf{positive} ($+$); heights downwards are \textbf{negative} ($-$).
\end{itemize}
\textbf{2. Reflection at Spherical Mirrors:}
\begin{equation}
  \mathbf{\frac{1}{v} + \frac{1}{u} = \frac{1}{f} = \frac{2}{R}}
\end{equation}
\begin{itemize}
  \item \textbf{Concave Mirror:} $f < 0, R < 0$ (converging). Produces real inverted images for $u > |f|$, and virtual erect magnified image for $u < |f|$.
  \item \textbf{Convex Mirror:} $f > 0, R > 0$ (diverging). Produces virtual, erect, diminished images for all real object locations.
  \item \textbf{Transverse (Linear) Magnification:}
  \begin{equation}
    \mathbf{m = \frac{h_i}{h_o} = -\frac{v}{u} = \frac{f}{f - u} = \frac{f - v}{f}}
  \end{equation}
  $m < 0 \implies$ real, inverted; $m > 0 \implies$ virtual, erect.
  \item \textbf{Longitudinal Magnification for Small Objects ($du$):} $m_L = \frac{dv}{du} = -m^2 = -\left(\frac{v}{u}\right)^2$.
\end{itemize}
\end{conceptbox}

\begin{formulabox}[Refraction at Plane Interfaces, Apparent Depth \& Total Internal Reflection]
\textbf{1. Snell's Law of Refraction:}
\begin{equation}
  \mathbf{n_1 \sin i = n_2 \sin r \implies \frac{\sin i}{\sin r} = \frac{n_2}{n_1} = \,^1n_2 = \frac{v_1}{v_2} = \frac{\lambda_1}{\lambda_2}}
\end{equation}
Frequency $\nu$ remains invariant across any refractive boundary: $\nu_1 = \nu_2 = \nu$.
\textbf{2. Apparent Depth \& Normal Shift in Dense Media:}
An object placed at real depth $d$ in a medium of refractive index $\mu$ viewed normally from air:
\begin{equation}
  \mathbf{d_{\text{apparent}} = \frac{d}{\mu}}, \qquad \mathbf{\text{Normal Shift } \Delta t = d - d_{\text{app}} = d\left(1 - \frac{1}{\mu}\right)}
\end{equation}
For composite multiple horizontal layers: $d_{\text{apparent, total}} = \sum_{i=1}^n \frac{d_i}{\mu_i}$.
\textbf{3. Total Internal Reflection (TIR):}
When light propagates from an optically denser medium ($n_1$) toward a rarer medium ($n_2 < n_1$):
\begin{itemize}
  \item Critical angle $i_c$: angle of incidence where angle of refraction $r = 90^\circ$:
  \begin{equation}
    n_1 \sin i_c = n_2 \sin 90^\circ \implies \mathbf{\sin i_c = \frac{n_2}{n_1} = \frac{1}{\mu}} \quad (\text{if rarer medium is air, } n_2 = 1)
  \end{equation}
  \item For $i > i_c$, $100\%$ of incident light energy reflects back into the denser medium with zero transmission loss.
  \item \textbf{Optical Fibres:} Comprise a central cylindrical core ($\mu_1$) surrounded by a cladding ($\mu_2 < \mu_1$). Light admitted within the cone of acceptance undergoes repeated total internal reflections.
  \begin{equation}
    \mathbf{\text{Numerical Aperture (NA)} = \sin\theta_a = \sqrt{\mu_1^2 - \mu_2^2}}
  \end{equation}
\end{itemize}
\end{formulabox}

\begin{formulabox}[Refraction at a Single Spherical Surface: First-Principles Derivation]
Consider a convex spherical refracting interface of radius of curvature $R$ separating medium 1 ($\mu_1$) from medium 2 ($\mu_2$). A point object $O$ lies on the principal axis in medium 1 at distance $-u$. Light refracts to form an image $I$ at $+v$ in medium 2.
\begin{enumerate}
  \item Let ray from $O$ strike surface at $M$ near pole $P$. Let angles with principal axis be $\alpha$ (at $O$), $\beta$ (at $I$), and $\theta$ (at center of curvature $C$).
  \item In triangle $OMC$, exterior angle is incidence angle: $i = \alpha + \theta$.
  \item In triangle $IMC$, exterior angle is $\theta$: $\theta = r + \beta \implies r = \theta - \beta$.
  \item Under paraxial approximations ($\tan\phi \approx \phi \approx \sin\phi$ and $M \approx P$):
  \begin{equation}
    \alpha \approx \frac{PM}{u} = -\frac{h}{u}, \qquad \beta \approx \frac{PM}{v} = \frac{h}{v}, \qquad \theta \approx \frac{PM}{R} = \frac{h}{R}
  \end{equation}
  \item By Snell's Law for small angles: $\mu_1 i = \mu_2 r$.
  \begin{equation}
    \mu_1 (\alpha + \theta) = \mu_2 (\theta - \beta) \implies \mu_1 \alpha + \mu_2 \beta = (\mu_2 - \mu_1)\theta
  \end{equation}
  \item Substituting the paraxial arc angles:
  \begin{equation}
    \mu_1 \left(-\frac{h}{u}\right) + \mu_2 \left(\frac{h}{v}\right) = (\mu_2 - \mu_1)\frac{h}{R}
  \end{equation}
  Cancelling $h$:
  \begin{equation}
    \mathbf{\frac{\mu_2}{v} - \frac{\mu_1}{u} = \frac{\mu_2 - \mu_1}{R}}
  \end{equation}
\end{enumerate}
This master governing formula holds universally for both convex and concave surfaces, whether light travels from rarer to denser or denser to rarer, provided proper Cartesian signs are used for $u, v, R$.
\end{formulabox}

\section{Thin Lenses, Lens-Maker's Formula, Prisms \& Optical Instruments}

\begin{formulabox}[Lens-Maker's Formula \& Thin Lens Equation: Step-by-Step Proof]
Consider a thin lens of refractive index $\mu_2$ bounded by spherical surfaces of radii $R_1$ and $R_2$ immersed in a surrounding medium of refractive index $\mu_1$.
\begin{enumerate}
  \item \textbf{Refraction at First Surface ($R_1$):} Object at $u$, forms intermediate image at $v'$:
  \begin{equation}
    \frac{\mu_2}{v'} - \frac{\mu_1}{u} = \frac{\mu_2 - \mu_1}{R_1} \tag{Eq. 1}
  \end{equation}
  \item \textbf{Refraction at Second Surface ($R_2$):} The intermediate image $v'$ serves as a virtual object for the second interface, forming final image at $v$:
  \begin{equation}
    \frac{\mu_1}{v} - \frac{\mu_2}{v'} = \frac{\mu_1 - \mu_2}{R_2} = -\frac{\mu_2 - \mu_1}{R_2} \tag{Eq. 2}
  \end{equation}
  \item \textbf{Adding Equations (1) and (2):} The intermediate term $\frac{\mu_2}{v'}$ cancels completely:
  \begin{equation}
    \mu_1 \left(\frac{1}{v} - \frac{1}{u}\right) = (\mu_2 - \mu_1)\left(\frac{1}{R_1} - \frac{1}{R_2}\right)
  \end{equation}
  Dividing through by $\mu_1$:
  \begin{equation}
    \frac{1}{v} - \frac{1}{u} = \left(\frac{\mu_2}{\mu_1} - 1\right)\left(\frac{1}{R_1} - \frac{1}{R_2}\right)
  \end{equation}
  \item Setting $u \to -\infty \implies v = f$:
  \begin{equation}
    \mathbf{\frac{1}{f} = \left(\frac{\mu_{\text{lens}}}{\mu_{\text{medium}}} - 1\right)\left(\frac{1}{R_1} - \frac{1}{R_2}\right)} \quad (\text{Lens-Maker's Formula})
  \end{equation}
  \begin{equation}
    \mathbf{\frac{1}{v} - \frac{1}{u} = \frac{1}{f}} \quad (\text{Gaussian Thin Lens Formula})
  \end{equation}
\end{enumerate}
\textbf{3. Optical Power ($P$):}
\begin{equation}
  \mathbf{P = \frac{1}{f\,[\mathrm{m}]}} \quad [\mathrm{Dioptres}, \mathrm{D} = \mathrm{m}^{-1}]
\end{equation}
\textbf{4. Combination of Thin Lenses:}
\begin{itemize}
  \item \textbf{Lenses in Contact:} $\mathbf{P_{\text{eq}} = P_1 + P_2 \implies \frac{1}{F_{\text{eq}}} = \frac{1}{f_1} + \frac{1}{f_2}}$, magnification $m = m_1 \times m_2$.
  \item \textbf{Separated by Distance $d$:} $\mathbf{P_{\text{eq}} = P_1 + P_2 - d P_1 P_2 \implies \frac{1}{F_{\text{eq}}} = \frac{1}{f_1} + \frac{1}{f_2} - \frac{d}{f_1 f_2}}$.
\end{itemize}
\end{formulabox}

\begin{formulabox}[Refraction Through Prisms, Minimum Deviation \& Dispersion]
A triangular prism of refracting angle $A$ and refractive index $\mu$:
\begin{enumerate}
  \item \textbf{Geometric Relations:}
  \begin{equation}
    r_1 + r_2 = A, \qquad \mathbf{\delta = i + e - A}
  \end{equation}
  where $i$ is incidence angle, $e$ is emergence angle, and $\delta$ is angular deviation.
  \item \textbf{Condition for Minimum Deviation ($\delta = \delta_m$):}
  By symmetry and reversibility of light paths: $i = e$ and $r_1 = r_2 = \frac{A}{2}$.
  \begin{equation}
    \delta_m = 2i - A \implies i = \frac{A + \delta_m}{2}
  \end{equation}
  Applying Snell's law at first refracting surface ($\sin i = \mu \sin r_1$):
  \begin{equation}
    \mathbf{\mu = \frac{\sin\left(\frac{A + \delta_m}{2}\right)}{\sin\left(\frac{A}{2}\right)}} \quad (\text{Prism Formula})
  \end{equation}
  \item \textbf{Thin Prism Approximation ($A \le 10^\circ$):}
  \begin{equation}
    \sin\theta \approx \theta \implies \frac{A + \delta}{2} \approx \mu \frac{A}{2} \implies \mathbf{\delta = (\mu - 1) A}
  \end{equation}
  \item \textbf{Dispersion \& Dispersive Power ($\omega$):}
  \begin{itemize}
    \item Angular Dispersion: $\mathbf{\theta = \delta_V - \delta_R = (\mu_V - \mu_R) A}$.
    \item Mean Deviation (Yellow light): $\delta_Y = (\mu_Y - 1) A$.
    \item Dispersive Power: $\mathbf{\omega = \frac{\delta_V - \delta_R}{\delta_Y} = \frac{\mu_V - \mu_R}{\mu_Y - 1}}$ (Independent of prism angle $A$).
  \end{itemize}
\end{enumerate}
\end{formulabox}

\begin{formulabox}[Magnifying Power of Optical Instruments: Microscopes \& Telescopes]
\textbf{1. Simple Microscope (Magnifying Glass):}
\begin{itemize}
  \item Image formed at Near Point ($D = 25\,\mathrm{cm}$): $\mathbf{m = 1 + \frac{D}{f}}$ (Maximum magnification, strained eye).
  \item Image formed at Infinity (Normal Adjustment): $\mathbf{m = \frac{D}{f}}$ (Relaxed eye).
\end{itemize}
\textbf{2. Compound Microscope:}
Objective lens ($f_o$, small aperture) + Eyepiece ($f_e$, larger aperture).
\begin{itemize}
  \item Total Magnifying Power: $m = m_o \times m_e = \left(-\frac{v_o}{u_o}\right) m_e$.
  \item \textbf{Near Point Adjustment ($D$):} $\mathbf{m = -\frac{v_o}{u_o}\left(1 + \frac{D}{f_e}\right)}$, Tube length $L = v_o + |u_e|$.
  \item \textbf{Normal Adjustment (Infinity):} $\mathbf{m \approx -\frac{L}{f_o}\frac{D}{f_e}}$, Tube length $L = v_o + f_e$.
\end{itemize}
\textbf{3. Astronomical Refracting Telescope:}
Objective ($f_o$, large aperture) + Eyepiece ($f_e$, small aperture).
\begin{itemize}
  \item \textbf{Normal Adjustment (Image at Infinity):} $\mathbf{m = -\frac{f_o}{f_e}}$, Tube length $\mathbf{L = f_o + f_e}$.
  \item \textbf{Near Point Adjustment ($D$):} $\mathbf{m = -\frac{f_o}{f_e}\left(1 + \frac{f_e}{D}\right)}$, Tube length $L = f_o + \frac{f_e D}{f_e + D}$.
\end{itemize}
\end{formulabox}

\begin{trapbox}[Syllabus Boundary Alert: Deletion of Resolving Power of Optical Instruments]
\textbf{NMC REVISION WARNING FOR NEET UG:}
Under the revised NEET UG curriculum:
\begin{itemize}
  \item \textbf{STRICTLY DELETED:} The theoretical formulas for \textbf{Resolving Power of a Microscope} ($\text{RP} = \frac{2\mu\sin\theta}{1.22\lambda}$) and \textbf{Resolving Power of an Astronomical Telescope} ($\text{RP} = \frac{D}{1.22\lambda}$), along with Rayleigh's criterion angular limit $\theta = \frac{1.22\lambda}{D}$, have been \textbf{completely deleted}.
  \item \textbf{EXAM ORIENTATION:} Candidates must focus on Magnifying Power derivations, tube lengths, and ray diagrams of microscopes/telescopes without computing resolving powers.
\end{itemize}
\end{trapbox}

\section{Wave Optics: Huygens' Principle, Interference \& Diffraction}

\begin{conceptbox}[Wavefronts, Huygens' Secondary Wavelets \& Interference Physics]
\textbf{1. Wavefront Classification:}
A wavefront is the locus of all adjacent points in a medium vibrating with identical instantaneous phase:
\begin{itemize}
  \item \textbf{Spherical Wavefront:} Emitted by an isotropic point source; intensity falls as $I \propto \frac{1}{r^2}$, amplitude $A \propto \frac{1}{r}$.
  \item \textbf{Cylindrical Wavefront:} Emitted by a linear slit source; intensity $I \propto \frac{1}{r}$, amplitude $A \propto \frac{1}{\sqrt{r}}$.
  \item \textbf{Plane Wavefront:} Formed at infinite distance from any source; intensity and amplitude remain constant.
  \item Light rays are everywhere perpendicular (normal) to the advancing wavefront.
\end{itemize}
\textbf{2. Huygens' Principle:}
\begin{enumerate}
  \item Every point on a primary wavefront acts as a source of secondary spherical wavelets spreading in all forward directions at the speed of the wave.
  \item The forward envelope (tangential surface) touching all these secondary wavelets at any subsequent instant defines the new position of the wavefront.
\end{enumerate}
\textbf{3. Superposition of Two Coherent Waves:}
Let two waves of identical frequency have amplitudes $a_1, a_2$ and relative phase difference $\phi$:
\begin{equation}
  y_1 = a_1 \sin(\omega t), \qquad y_2 = a_2 \sin(\omega t + \phi)
\end{equation}
Resultant amplitude $A_R$ and resultant intensity $I_R$:
\begin{equation}
  \mathbf{A_R = \sqrt{a_1^2 + a_2^2 + 2a_1 a_2 \cos\phi}}, \qquad \mathbf{I_R = I_1 + I_2 + 2\sqrt{I_1 I_2}\cos\phi}
\end{equation}
If two sources are identical ($I_1 = I_2 = I_0$):
\begin{equation}
  \mathbf{I_R = 2I_0(1 + \cos\phi) = 4I_0 \cos^2\left(\frac{\phi}{2}\right)}
\end{equation}
\begin{itemize}
  \item \textbf{Constructive Interference (Bright Fringes):} $\phi = 2n\pi \iff \Delta x = n\lambda \implies \mathbf{I_{\max} = (\sqrt{I_1} + \sqrt{I_2})^2 = 4I_0}$.
  \item \textbf{Destructive Interference (Dark Fringes):} $\phi = (2n-1)\pi \iff \Delta x = (2n-1)\frac{\lambda}{2} \implies \mathbf{I_{\min} = (\sqrt{I_1} - \sqrt{I_2})^2 = 0}$.
  \item \textbf{Fringe Contrast / Visibility:} $V = \frac{I_{\max} - I_{\min}}{I_{\max} + I_{\min}} = \frac{2\sqrt{I_1 I_2}}{I_1 + I_2}$.
\end{itemize}
\end{conceptbox}

\begin{formulabox}[Young's Double Slit Experiment (YDSE): Path Difference \& Fringe Width]
Two narrow slits $S_1, S_2$ separated by distance $d$ illuminate a screen at distance $D \gg d$. Let $P$ be a point at height $y$ from the central line.
\begin{enumerate}
  \item Path difference $\Delta x = S_2 P - S_1 P$:
  \begin{equation}
    S_2 P^2 - S_1 P^2 = \left[D^2 + \left(y + \frac{d}{2}\right)^2\right] - \left[D^2 + \left(y - \frac{d}{2}\right)^2\right] = 2yd
  \end{equation}
  \begin{equation}
    (S_2 P - S_1 P)(S_2 P + S_1 P) = 2yd \implies \Delta x \approx \frac{2yd}{2D} \implies \mathbf{\Delta x = \frac{yd}{D} = d\sin\theta}
  \end{equation}
  \item \textbf{Bright Fringe Positions (Maxima):}
  \begin{equation}
    \Delta x = n\lambda \implies \frac{y_n d}{D} = n\lambda \implies \mathbf{y_n = \frac{n\lambda D}{d}} \quad (n = 0, \pm 1, \pm 2, \dots)
  \end{equation}
  \item \textbf{Dark Fringe Positions (Minima):}
  \begin{equation}
    \Delta x = (2n-1)\frac{\lambda}{2} \implies \mathbf{y_n' = (2n-1)\frac{\lambda D}{2d}} \quad (n = 1, 2, 3, \dots)
  \end{equation}
  \item \textbf{Fringe Width ($\beta$):} Distance separating any two consecutive bright or dark fringes:
  \begin{equation}
    \mathbf{\beta = y_{n+1} - y_n = \frac{\lambda D}{d}}
  \end{equation}
  $\beta$ is completely uniform across all fringe orders on the screen.
  \item \textbf{Angular Fringe Width ($\theta_0$):}
  \begin{equation}
    \mathbf{\theta_0 = \frac{\beta}{D} = \frac{\lambda}{d}} \quad (\text{Independent of screen distance } D)
  \end{equation}
  \item \textbf{Fringe Shift due to Transparent Thin Sheet of Thickness $t$ and Index $\mu$:}
  Introducing the sheet in front of one slit introduces additional optical path $(\mu - 1)t$:
  \begin{equation}
    \mathbf{\Delta y_0 = \frac{(\mu - 1)t D}{d} = \frac{(\mu - 1)t}{\lambda}\beta}
  \end{equation}
  The entire fringe pattern shifts toward the covered slit without changing fringe width $\beta$.
\end{enumerate}
\end{formulabox}

\begin{formulabox}[Single-Slit Fraunhofer Diffraction: Minima \& Central Maximum Width]
Light of wavelength $\lambda$ illuminates a single rectangular slit of width $a$.
\begin{enumerate}
  \item \textbf{Diffraction Minima Condition:}
  The slit aperture is divided into $2n$ equal zones that cancel pairwise:
  \begin{equation}
    \mathbf{a \sin\theta = n \lambda} \quad (n = \pm 1, \pm 2, \pm 3, \dots)
  \end{equation}
  For small angles: $\theta_n \approx \frac{n\lambda}{a}$.
  \item \textbf{Secondary Maxima Condition:}
  \begin{equation}
    \mathbf{a \sin\theta \approx (2n+1)\frac{\lambda}{2}} \quad (n = 1, 2, 3, \dots)
  \end{equation}
  Peak intensities drop precipitously: $I_0 : \frac{I_0}{22} : \frac{I_0}{61} : \dots$
  \item \textbf{Angular Width of Central Maximum ($2\theta_0$):}
  Bounded by the first diffraction minima on either side ($n = \pm 1$):
  \begin{equation}
    \mathbf{2\theta_0 = \frac{2\lambda}{a}}
  \end{equation}
  \item \textbf{Linear Width of Central Diffraction Maximum ($W$):}
  \begin{equation}
    \mathbf{W = 2 y_1 = 2 \left(\frac{\lambda D}{a}\right) = \frac{2\lambda D}{a}}
  \end{equation}
  Notice: The central diffraction peak is \textbf{twice as wide} as secondary maxima ($W = 2\beta_{\text{diff}}$).
\end{enumerate}
\end{formulabox}

\begin{trapbox}[Syllabus Boundary Alert: Complete Deletion of Light Polarisation]
\textbf{NMC REVISION WARNING FOR NEET UG:}
Under the revised NEET UG syllabus issued by the National Medical Commission:
\begin{itemize}
  \item \textbf{STRICTLY DELETED:} The entire topic of \textbf{Polarisation of Light} has been \textbf{completely eliminated}.
  \item Deleted concepts include: Brewster's Law of polarisation by reflection ($\tan i_p = \mu$), Brewster's angle, Malus' Law of transmitted intensity ($I = I_0 \cos^2\theta$), Polaroid sheets, polarisation by scattering, and optical activity.
  \item \textbf{EXAM ORIENTATION:} Candidates must focus solely on Wavefronts, Huygens' Principle, YDSE Interference, and Single-Slit Fraunhofer Diffraction.
\end{itemize}
\end{trapbox}

\section{Dual Nature of Matter \& Radiation: Photoelectric Effect \& de Broglie Waves}

\begin{conceptbox}[Photoelectric Effect: Experimental Laws \& Einstein's Quantum Derivation]
\textbf{1. Key Experimental Observations (Hertz, Hallwachs, Lenard):}
\begin{itemize}
  \item \textbf{Instantaneous Emission:} Electron emission occurs without any measurable time lag ($< 10^{-9}\,\mathrm{s}$), contradicting classical wave continuous energy buildup.
  \item \textbf{Threshold Frequency ($\nu_0$):} For each photosensitive metal, there exists a minimum characteristic threshold frequency $\nu_0$ below which \textbf{no photoelectrons are emitted}, regardless of incident light intensity or duration.
  \item \textbf{Kinetic Energy vs. Frequency:} Maximum kinetic energy $K_{\max}$ is strictly linear with incident frequency $\nu$ and completely independent of light intensity.
  \item \textbf{Saturation Current vs. Intensity:} Photoelectric saturation current is strictly proportional to incident light intensity (photon flux).
\end{itemize}
\textbf{2. Einstein's Photoelectric Equation Derivation:}
Light consists of discrete energy quanta called \textbf{photons}, each carrying energy $E = h\nu$.
When a photon of energy $h\nu$ collides with a conduction electron, its entire energy is transferred to that single electron. A portion equal to the metal's \textbf{work function} ($\Phi_0 = h\nu_0$) is consumed to overcome surface binding; the remainder appears as maximum electron kinetic energy:
\begin{equation}
  h\nu = \Phi_0 + K_{\max} \implies \mathbf{K_{\max} = h\nu - \Phi_0 = h(\nu - \nu_0) = \frac{hc}{\lambda} - \frac{hc}{\lambda_0}}
\end{equation}
where $h = 6.626 \times 10^{-34}\,\mathrm{J}\cdot\mathrm{s}$ is Planck's constant ($hc \approx 12400\,\mathrm{eV}\cdot\text{\AA} \approx 1240\,\mathrm{eV}\cdot\mathrm{nm}$).
\textbf{3. Stopping Potential ($V_0$):}
The negative retarding collector potential required to halt the most energetic photoelectrons:
\begin{equation}
  K_{\max} = e V_0 \implies \mathbf{V_0 = \left(\frac{h}{e}\right)\nu - \frac{\Phi_0}{e} = \left(\frac{h}{e}\right)(\nu - \nu_0)}
\end{equation}
\textit{Universal Graph Property:} In the $V_0$ vs. $\nu$ straight-line graph, the slope is a universal constant:
\begin{equation}
  \mathbf{\text{Slope} = \frac{h}{e} \approx 4.14 \times 10^{-15}\,\mathrm{V}\cdot\mathrm{s}}
\end{equation}
The horizontal axis intercept is threshold frequency $\nu_0$; vertical axis intercept is $-\Phi_0/e$.
\end{conceptbox}

\begin{formulabox}[de Broglie Matter Waves: Particle-Wave Duality Expressions]
Louis de Broglie postulated that any moving material particle of mass $m$ and velocity $v$ possesses an associated matter wave:
\begin{equation}
  \mathbf{\lambda = \frac{h}{p} = \frac{h}{m v} = \frac{h}{\sqrt{2 m K}}}
\end{equation}
where $K$ is kinetic energy.
\textbf{1. De Broglie Wavelength of Charged Particles Accelerated through Potential $V$:}
$K = q V \implies \lambda = \frac{h}{\sqrt{2 m q V}}$.
\begin{itemize}
  \item \textbf{Electron ($m = 9.109 \times 10^{-31}\,\mathrm{kg}, q = e$):}
  \begin{equation}
    \mathbf{\lambda_e = \frac{1.227}{\sqrt{V}}\,\mathrm{nm} = \frac{12.27}{\sqrt{V}}\,\text{\AA} = \sqrt{\frac{150}{V}}\,\text{\AA}}
  \end{equation}
  \item \textbf{Proton ($m_p = 1.673 \times 10^{-27}\,\mathrm{kg}, q = e$):}
  \begin{equation}
    \mathbf{\lambda_p = \frac{0.286}{\sqrt{V}}\,\text{\AA}}
  \end{equation}
  \item \textbf{Deuteron ($m_d = 2m_p, q = e$):}
  \begin{equation}
    \mathbf{\lambda_d = \frac{0.202}{\sqrt{V}}\,\text{\AA}}
  \end{equation}
  \item \textbf{Alpha Particle ($m_\alpha = 4m_p, q = 2e$):}
  \begin{equation}
    \mathbf{\lambda_\alpha = \frac{0.101}{\sqrt{V}}\,\text{\AA}}
  \end{equation}
\end{itemize}
\textbf{2. Thermal Neutrons in Equilibrium at Temperature $T$:}
Mean translational kinetic energy $K = \frac{3}{2}k_B T$:
\begin{equation}
  \mathbf{\lambda_{\text{thermal}} = \frac{h}{\sqrt{3 m k_B T}}}
\end{equation}
\end{formulabox}

\begin{trapbox}[Syllabus Boundary Alert: Deletion of Davisson-Germer Experiment]
\textbf{NMC REVISION WARNING FOR NEET UG:}
Under the revised NEET UG curriculum:
\begin{itemize}
  \item \textbf{STRICTLY DELETED:} The \textbf{Davisson-Germer Experiment} (nickel crystal electron diffraction, Bragg peak at $54\,\mathrm{V}$ and $\theta = 50^\circ$, accelerating electron gun details) has been \textbf{formally removed}.
  \item \textbf{EXAM ORIENTATION:} Candidates must be proficient in Einstein's photoelectric equation, photon momentum calculations, and de Broglie wavelength formulas for accelerated particles without studying the experimental apparatus of Davisson and Germer.
\end{itemize}
\end{trapbox}

\section{Atomic Physics: Bohr's Postulates, Spectral Series \& Energy Levels}

\begin{conceptbox}[Rutherford Scattering \& Bohr's Atomic Postulates]
\textbf{1. Rutherford Scattering (Geiger-Marsden Experiment):}
Distance of closest approach $r_0$ for an $\alpha$-particle ($q_\alpha = 2e$) directed head-on with kinetic energy $K$ toward a nucleus of atomic number $Z$:
\begin{equation}
  K = \frac{1}{4\pi\varepsilon_0}\frac{(2e)(Ze)}{r_0} \implies \mathbf{r_0 = \frac{1}{4\pi\varepsilon_0}\frac{2 Z e^2}{K}}
\end{equation}
Impact parameter: $b = \frac{1}{4\pi\varepsilon_0}\frac{Z e^2 \cot(\theta/2)}{K}$.
\textbf{2. Bohr's Foundational Postulates for Hydrogen-Like Species:}
\begin{enumerate}
  \item \textbf{Stable Stationary Orbits:} An orbital electron circles the positive nucleus solely under the Coulomb attraction without radiating electromagnetic energy, despite continuous centripetal acceleration.
  \item \textbf{Quantization of Orbital Angular Momentum:} Stable orbits are restricted strictly to those where orbital angular momentum is an integer multiple of $\hbar = h / 2\pi$:
  \begin{equation}
    \mathbf{L = m_e v r = \frac{n h}{2\pi}} \quad (n = 1, 2, 3, \dots)
  \end{equation}
  \item \textbf{Frequency Condition (Radiative Transitions):} Photons are emitted or absorbed exclusively when an electron undergoes a discrete transition between stationary energy levels $E_i$ and $E_f$:
  \begin{equation}
    \mathbf{\Delta E = E_i - E_f = h\nu = \frac{hc}{\lambda}}
  \end{equation}
\end{enumerate}
\end{conceptbox}

\begin{formulabox}[Bohr Atom Formulas: Rigorous Analytical Derivations for Hydrogen-Like Species]
For a hydrogen-like atom of nuclear charge $+Ze$ and electron mass $m$:
\begin{enumerate}
  \item \textbf{Electrostatic Centripetal Equilibrium:}
  \begin{equation}
    \frac{m v^2}{r} = \frac{1}{4\pi\varepsilon_0} \frac{Z e^2}{r^2} \implies m v^2 r = \frac{Z e^2}{4\pi\varepsilon_0} \tag{Eq. 1}
  \end{equation}
  \item \textbf{Angular Momentum Quantization:}
  \begin{equation}
    m v r = \frac{n h}{2\pi} \implies v = \frac{n h}{2\pi m r} \tag{Eq. 2}
  \end{equation}
  \item \textbf{Orbital Radius Derivation:} Substituting Eq. 2 into Eq. 1:
  \begin{equation}
    m \left(\frac{n^2 h^2}{4\pi^2 m^2 r^2}\right) r = \frac{Z e^2}{4\pi\varepsilon_0} \implies \mathbf{r_n = \frac{\varepsilon_0 n^2 h^2}{\pi m Z e^2} = r_0 \frac{n^2}{Z} = 0.529 \frac{n^2}{Z}\,\text{\AA}}
  \end{equation}
  where $r_0 = 0.529177\,\text{\AA}$ is the first Bohr radius of hydrogen ($n=1, Z=1$).
  \item \textbf{Orbital Velocity Derivation:}
  \begin{equation}
    \mathbf{v_n = \frac{Z e^2}{2\varepsilon_0 n h} = v_0 \frac{Z}{n} = \left(2.18 \times 10^6\,\mathrm{m/s}\right)\frac{Z}{n} = c \alpha \frac{Z}{n}}
  \end{equation}
  where $\alpha = \frac{e^2}{2\varepsilon_0 h c} \approx \frac{1}{137}$ is the fine-structure constant.
  \item \textbf{Energies in the $n$-th Stationary Orbit:}
  \begin{itemize}
    \item \textbf{Kinetic Energy:} $\mathbf{K_n = \frac{1}{2}m v_n^2 = \frac{m Z^2 e^4}{8\varepsilon_0^2 n^2 h^2} = +13.6 \frac{Z^2}{n^2}\,\mathrm{eV}}$.
    \item \textbf{Potential Energy:} $\mathbf{U_n = -\frac{1}{4\pi\varepsilon_0}\frac{Z e^2}{r_n} = -\frac{m Z^2 e^4}{4\varepsilon_0^2 n^2 h^2} = -27.2 \frac{Z^2}{n^2}\,\mathrm{eV}}$.
    \item \textbf{Total Mechanical Energy:}
    \begin{equation}
      \mathbf{E_n = K_n + U_n = -\frac{m Z^2 e^4}{8\varepsilon_0^2 n^2 h^2} = -13.6 \frac{Z^2}{n^2}\,\mathrm{eV}}
    \end{equation}
    \item \textbf{Golden Energy Ratios:}
    \begin{equation}
      \mathbf{E_n = -K_n = \frac{U_n}{2}}
    \end{equation}
  \end{itemize}
\end{enumerate}
\end{formulabox}

\begin{formulabox}[Hydrogen Spectral Emission Series \& Rydberg Equation]
When an electron drops from an upper stationary orbit $n_2$ to a lower orbit $n_1$ ($n_2 > n_1$):
\begin{equation}
  \Delta E = E_{n_2} - E_{n_1} = 13.6 Z^2 \left(\frac{1}{n_1^2} - \frac{1}{n_2^2}\right)\,\mathrm{eV}
\end{equation}
\begin{equation}
  \mathbf{\frac{1}{\lambda} = \overline{\nu} = R_\infty Z^2 \left(\frac{1}{n_1^2} - \frac{1}{n_2^2}\right)}
\end{equation}
where $R_\infty = \frac{m e^4}{8\varepsilon_0^2 c h^3} \approx 1.09737 \times 10^7\,\mathrm{m}^{-1}$ is the Rydberg constant ($\frac{1}{R_\infty} \approx 912\,\text{\AA}$).
\begin{center}
\begin{tabularx}{\linewidth}{|l|c|c|c|Y|}
\hline
\textbf{Spectral Series} & \textbf{Lower Level ($n_1$)} & \textbf{Upper Level ($n_2$)} & \textbf{Spectral Domain} & \textbf{Wavelength Range} \\
\hline
\textbf{Lyman} & $1$ & $2, 3, 4, \dots$ & \textbf{Ultraviolet (UV)} & $\lambda_{\min} = 91.2\,\mathrm{nm}, \lambda_{\max} = 121.6\,\mathrm{nm}$ \\
\hline
\textbf{Balmer} & $2$ & $3, 4, 5, \dots$ & \textbf{Visible} (part UV) & $\lambda_{\min} = 364.6\,\mathrm{nm}, \lambda_{\max} = 656.3\,\mathrm{nm}$ \\
\hline
\textbf{Paschen} & $3$ & $4, 5, 6, \dots$ & \textbf{Near Infrared} & $\lambda_{\min} = 820.4\,\mathrm{nm}, \lambda_{\max} = 1875\,\mathrm{nm}$ \\
\hline
\textbf{Brackett} & $4$ & $5, 6, 7, \dots$ & \textbf{Mid Infrared} & $\lambda_{\min} = 1458\,\mathrm{nm}, \lambda_{\max} = 4051\,\mathrm{nm}$ \\
\hline
\textbf{Pfund} & $5$ & $6, 7, 8, \dots$ & \textbf{Far Infrared} & $\lambda_{\min} = 2279\,\mathrm{nm}, \lambda_{\max} = 7458\,\mathrm{nm}$ \\
\hline
\end{tabularx}
\end{center}
\textbf{Total Spectral Lines Emitted from Excited Level $n$ to Ground State:}
\begin{equation}
  \mathbf{N = \frac{n(n - 1)}{2}}
\end{equation}
\end{formulabox}

\section{Nuclear Physics: Composition, Density, Mass Defect \& Binding Energy}

\begin{conceptbox}[Nuclear Dimensions, Universal Nuclear Density \& Strong Nuclear Force]
\textbf{1. Nuclear Constituents:}
Atomic nucleus contains $Z$ protons and $N = A - Z$ neutrons (nucleons).
\textbf{2. Nuclear Radius Empirical Law:}
\begin{equation}
  \mathbf{R = R_0 A^{1/3}}
\end{equation}
where $R_0 \approx 1.2 \times 10^{-15}\,\mathrm{m} = 1.2\,\mathrm{fm}$.
\textbf{3. Proof of Constant Nuclear Density ($\rho_{\text{nuc}}$):}
Nuclear volume: $V = \frac{4}{3}\pi R^3 = \frac{4}{3}\pi (R_0 A^{1/3})^3 = \frac{4}{3}\pi R_0^3 A$.
Nuclear mass: $M \approx A \cdot m_N$ (where nucleon mass $m_N \approx 1.66 \times 10^{-27}\,\mathrm{kg}$).
\begin{equation}
  \mathbf{\rho_{\text{nuc}} = \frac{M}{V} = \frac{A m_N}{\frac{4}{3}\pi R_0^3 A} = \frac{3 m_N}{4\pi R_0^3} \approx 2.3 \times 10^{17}\,\mathrm{kg/m}^3}
\end{equation}
Nuclear density is an \textbf{absolute universal constant}, completely independent of mass number $A$.
\textbf{4. Properties of the Strong Nuclear Force:}
\begin{itemize}
  \item Strongest fundamental force in nature ($\sim 100$ times stronger than electrostatic repulsion).
  \item Strictly short-ranged: active solely across $1 - 2\,\mathrm{fm}$. For $r < 0.8\,\mathrm{fm}$, it becomes strongly repulsive.
  \item Charge-independent: $F_{p-p} = F_{p-n} = F_{n-n}$ under identical quantum states.
  \item Saturated: each nucleon interacts exclusively with its immediate neighbouring nucleons.
\end{itemize}
\end{conceptbox}

\begin{formulabox}[Mass Defect, Nuclear Binding Energy Curve \& Energy Reactions]
\textbf{1. Mass Defect ($\Delta m$):}
The rest mass of any stable bound nucleus is strictly less than the combined rest masses of its constituent free protons and neutrons:
\begin{equation}
  \mathbf{\Delta m = [Z m_p + (A - Z) m_n] - M_{\text{nucleus}}(A, Z)}
\end{equation}
\textbf{2. Nuclear Binding Energy ($E_b$):}
By Einstein's mass-energy equivalence ($E = \Delta m c^2$):
\begin{equation}
  \mathbf{E_b = \Delta m c^2 = \Delta m\,[\mathrm{u}] \times 931.5\,\mathrm{MeV}}
\end{equation}
where $1\,\mathrm{u} \equiv 931.494\,\mathrm{MeV}/c^2$.
\textbf{3. Binding Energy per Nucleon ($E_b/A$ Curve):}
The stability of a nucleus is dictated strictly by binding energy per nucleon ($E_b/A$):
\begin{itemize}
  \item Rapid increase from light nuclei ($^2\mathrm{H} \approx 1.1\,\mathrm{MeV/nucleon}$) with local peaks at tightly bound alpha-multiples ($^4\mathrm{He}, ^{12}\mathrm{C}, ^{16}\mathrm{O}$).
  \item Broad flat maximum in intermediate mass range $30 < A < 120$ with peak at $\mathbf{^{56}\mathrm{Fe} \approx 8.75\,\mathrm{MeV/nucleon}}$.
  \item Gradual drop toward heavy nuclei due to cumulative Coulomb repulsion between protons, dropping to $\approx 7.6\,\mathrm{MeV/nucleon}$ for $^{238}\mathrm{U}$.
  \item \textbf{Nuclear Fission:} Heavy nucleus splitting ($A \sim 240 \to A_1 + A_2 \sim 120$) gains $\approx 0.9\,\mathrm{MeV/nucleon} \implies Q \approx 240 \times 0.9 \approx \mathbf{200\,\mathrm{MeV}}$ per fission.
  \item \textbf{Nuclear Fusion:} Two very light nuclei fusing into a heavier nucleus ($^2\mathrm{H} + ^3\mathrm{H} \to ^4\mathrm{He} + n$) releases $\approx \mathbf{17.6\,\mathrm{MeV}}$ ($3.5\,\mathrm{MeV/nucleon}$).
\end{itemize}
\end{formulabox}

\begin{trapbox}[Syllabus Boundary Alert: Complete Deletion of Radioactivity]
\textbf{NMC REVISION WARNING FOR NEET UG:}
Under the revised NEET UG syllabus issued by the National Medical Commission:
\begin{itemize}
  \item \textbf{STRICTLY DELETED:} The entire subtopic of \textbf{Radioactivity} has been \textbf{completely eliminated}.
  \item Deleted topics include: $\alpha$-decay, $\beta$-decay ($\beta^-$ and $\beta^+$), $\gamma$-decay, Radioactive Decay Law ($N(t) = N_0 e^{-\lambda t}$), Half-life ($T_{1/2} = \frac{0.693}{\lambda}$), Mean Life ($\tau = \frac{1}{\lambda}$), and Activity ($A = \lambda N$).
  \item \textbf{EXAM ORIENTATION:} Candidates must focus on Mass Defect, Binding Energy per Nucleon, Nuclear Density, and Fission/Fusion $Q$-values without solving radioactive decay kinematics.
\end{itemize}
\end{trapbox}

\section{Semiconductor Physics, Diodes \& Digital Logic Gates}

\begin{conceptbox}[Energy Bands, Intrinsic \& Extrinsic Semiconductors]
\textbf{1. Energy Band Theory of Solids:}
\begin{itemize}
  \item \textbf{Conductors:} Conduction band and valence band overlap, or conduction band is partially filled ($E_g = 0$). High electron density $n \sim 10^{28}\,\mathrm{m}^{-3}$.
  \item \textbf{Insulators:} Large forbidden energy gap: $\mathbf{E_g > 3.0\,\mathrm{eV}}$ (e.g., Diamond $E_g \approx 5.4\,\mathrm{eV}$).
  \item \textbf{Semiconductors:} Narrow forbidden energy gap: $\mathbf{E_g \le 3.0\,\mathrm{eV}}$ ($\mathrm{Si}: E_g = 1.1\,\mathrm{eV}$; $\mathrm{Ge}: E_g = 0.72\,\mathrm{eV}$; $\mathrm{GaAs}: E_g = 1.42\,\mathrm{eV}$).
\end{itemize}
\textbf{2. Intrinsic Semiconductors (Pure Crystals):}
Thermal energy excites valence electrons across $E_g$ into the conduction band, leaving equal numbers of vacant states (\textbf{holes}) in the valence band:
\begin{equation}
  \mathbf{n_e = n_h = n_i}
\end{equation}
Total electrical conductivity: $\mathbf{\sigma = e(n_e \mu_e + n_h \mu_h)}$.
\textbf{3. Extrinsic Semiconductors (Doped Crystals):}
\begin{itemize}
  \item \textbf{n-Type:} Doped with pentavalent donor impurities (P, As, Sb). Introduces donor discrete energy level $E_d$ just below conduction band ($0.05\,\mathrm{eV}$ for Si).
  \begin{equation}
    \mathbf{n_e \approx N_D \gg n_h} \quad (\text{Electrons: majority carriers; Holes: minority carriers})
  \end{equation}
  \item \textbf{p-Type:} Doped with trivalent acceptor impurities (B, Al, In, Ga). Introduces acceptor discrete energy level $E_a$ just above valence band ($0.05\,\mathrm{eV}$ for Si).
  \begin{equation}
    \mathbf{n_h \approx N_A \gg n_e} \quad (\text{Holes: majority carriers; Electrons: minority carriers})
  \end{equation}
  \item \textbf{Law of Mass Action (Thermal Equilibrium):}
  \begin{equation}
    \mathbf{n_e \cdot n_h = n_i^2(T)}
  \end{equation}
  Both n-type and p-type semiconductors are strictly \textbf{electrically neutral} overall.
\end{itemize}
\end{conceptbox}

\begin{formulabox}[p-n Junction Diode Physics, Biasing \& Rectifiers]
\textbf{1. p-n Junction Formation:}
Majority holes from p-side diffuse into n-side, and majority electrons from n-side diffuse into p-side.
Uncompensated immobile donor ions ($+$) on n-side and acceptor ions ($-$) on p-side form a \textbf{depletion layer} and built-in \textbf{barrier potential} ($V_0 \approx 0.7\,\mathrm{V}$ for Si, $0.3\,\mathrm{V}$ for Ge) pointing from n to p.
\textbf{2. Biasing Modes:}
\begin{itemize}
  \item \textbf{Forward Bias (p to $+$, n to $-$):} External field opposes barrier field $\implies$ barrier potential lowers to $(V_0 - V)$, depletion layer narrows, forward current rises exponentially: $I \approx I_0 (e^{eV/\eta k_B T} - 1)$.
  \item \textbf{Reverse Bias (p to $-$, n to $+$):} External field aids barrier field $\implies$ barrier increases to $(V_0 + V)$, depletion layer widens, current is tiny minority saturation current $I_0 \sim \mu\mathrm{A}$ (Si: $\mathrm{nA}$).
\end{itemize}
\textbf{3. Diode Rectifiers:}
\begin{itemize}
  \item \textbf{Half-Wave Rectifier:} One diode conducting during positive AC half-cycles.
  \begin{equation}
    \mathbf{\eta_{\max} = \frac{0.406 R_L}{r_f + R_L} \approx 40.6\%}, \qquad \mathbf{\nu_{\text{ripple}} = \nu_{\text{AC}}}
  \end{equation}
  \item \textbf{Full-Wave Rectifier (Center-Tapped or Bridge):} Two or four diodes conducting alternately.
  \begin{equation}
    \mathbf{\eta_{\max} = \frac{0.812 R_L}{r_f + R_L} \approx 81.2\%}, \qquad \mathbf{\nu_{\text{ripple}} = 2\nu_{\text{AC}}}
  \end{equation}
\end{itemize}
\end{formulabox}

\begin{formulabox}[Special Purpose Diodes: Zener, Photodiode, LED \& Solar Cell]
\textbf{1. Zener Diode as a Voltage Regulator:}
Heavily doped p-n junction with an ultra-thin depletion layer ($\sim 10^{-6}\,\mathrm{m}$) that undergoes sharp reverse breakdown at a fixed Zener breakdown voltage $V_z$.
Connected in \textbf{reverse bias} parallel with fluctuating load $R_L$:
\begin{equation}
  V_{\text{out}} = V_z = \text{constant}, \qquad I_{\text{supply}} = I_z + I_L
\end{equation}
Series current-limiting resistor $R_s$:
\begin{equation}
  \mathbf{R_s = \frac{V_{\text{unregulated}} - V_z}{I_z + I_L}}
\end{equation}
\textbf{2. Photodiode:} Operates in \textbf{reverse bias}. Incident photons ($h\nu > E_g$) generate electron-hole pairs in the depletion region. Fractional change in minority reverse current is much larger than in forward current, making reverse bias ideal for optical detection.
\textbf{3. Light Emitting Diode (LED):} Operates in \textbf{forward bias}. Electrons recombine with holes at junction, releasing recombination energy as spontaneous photons across direct bandgap: $\mathbf{\lambda = \frac{hc}{E_g}}$.
\textbf{4. Solar Cell:} Unbiased p-n junction operating in the \textbf{fourth quadrant} of $I-V$ characteristic, converting optical solar flux into electrical power.
\end{formulabox}

\begin{formulabox}[Digital Logic Gates: Truth Tables \& De Morgan's Theorems]
\textbf{1. Fundamental Logic Gates:}
\begin{itemize}
  \item \textbf{AND Gate:} $Y = A \cdot B$ (Output is 1 strictly when both inputs are 1).
  \item \textbf{OR Gate:} $Y = A + B$ (Output is 1 when at least one input is 1).
  \item \textbf{NOT Gate (Inverter):} $Y = \overline{A}$ ($0 \to 1, 1 \to 0$).
\end{itemize}
\textbf{2. Universal Logic Gates (Can Synthesize Any Boolean Function):}
\begin{itemize}
  \item \textbf{NAND Gate:} $Y = \overline{A \cdot B}$.
  \item \textbf{NOR Gate:} $Y = \overline{A + B}$.
\end{itemize}
\textbf{3. De Morgan's Master Theorems:}
\begin{equation}
  \mathbf{\overline{A + B} = \overline{A} \cdot \overline{B}}, \qquad \mathbf{\overline{A \cdot B} = \overline{A} + \overline{B}}
\end{equation}
\end{formulabox}

\begin{trapbox}[Syllabus Boundary Alert: Complete Deletion of Transistors]
\textbf{NMC REVISION WARNING FOR NEET UG:}
Under the revised NEET UG syllabus issued by the National Medical Commission:
\begin{itemize}
  \item \textbf{STRICTLY DELETED:} The entire topic of \textbf{Transistors} (bipolar junction transistors npn and pnp, emitter, base, collector actions) has been \textbf{completely eliminated}.
  \item Deleted subtopics include: Common Emitter (CE) input and output characteristics, current gain factors ($\alpha = I_c/I_e, \beta = I_c/I_b, \beta = \frac{\alpha}{1-\alpha}$), transistor as a CE amplifier (voltage gain $A_v = -\beta \frac{R_L}{R_i}$), transistor as a digital switch, and transistor as an oscillator.
  \item \textbf{EXAM ORIENTATION:} Candidates must focus exclusively on Semiconductors, p-n junction diodes, rectifiers, Zener diode regulators, optoelectronic diodes (LED, photodiode, solar cell), and Digital Logic Gates.
\end{itemize}
\end{trapbox}

\begin{ncertbox}[NCERT High-Yield Extraction: Optoelectronic Materials \& Eye Sensitivity]
The semiconductor bandgap for visible LEDs must exceed $1.8\,\mathrm{eV}$ ($hc/\lambda_{\max} \approx 1240 / 700\,\mathrm{nm} \approx 1.8\,\mathrm{eV}$).
\begin{itemize}
  \item \textbf{Gallium Arsenide Phosphide ($\mathrm{GaAs}_{1-x}\mathrm{P}_x$):} Used for red, amber, and yellow LEDs ($E_g \approx 1.9\,\mathrm{eV}$).
  \item \textbf{Gallium Nitride ($\mathrm{GaN}$):} Direct bandgap of $3.4\,\mathrm{eV}$ enables bright blue and ultraviolet LEDs.
  \item \textbf{Photodiode Detection Limit:} To detect an optical signal of wavelength $\lambda$, the semiconductor bandgap must satisfy $\mathbf{E_g \le \frac{hc}{\lambda}}$.
\end{itemize}
\end{ncertbox}

\section{Master High-Yield Problem Archetypes \& Worked Solutions}

\begin{usecasebox}[Archetype 1: Silvered Plano-Convex Lens Equivalent Mirror Focal Length]
\textbf{Problem:} A plano-convex lens of refractive index $\mu = 1.5$ has a curved surface of radius of curvature $R = 20\,\mathrm{cm}$. Compute the equivalent focal length when: (i) the flat plane surface is silvered, and (ii) the curved surface is silvered.
\textbf{Step-by-Step Solution:}
\begin{enumerate}
  \item Power of the plano-convex lens in air:
  \begin{equation}
    P_L = \frac{1}{f_L} = (\mu - 1)\left(\frac{1}{R_1} - \frac{1}{R_2}\right) = (1.5 - 1)\left(\frac{1}{20} - 0\right) = \frac{0.5}{20} = \frac{1}{40}\,\mathrm{cm}^{-1}
  \end{equation}
  \item Any silvered lens system acts as an equivalent mirror of power:
  \begin{equation}
    P_{\text{eq}} = 2 P_L + P_M \implies \mathbf{F_{\text{eq}} = -\frac{1}{P_{\text{eq}}}}
  \end{equation}
  \item \textbf{Case (i): Flat Plane Surface is Silvered:}
  Plane mirror has $R_M = \infty \implies P_M = -\frac{1}{f_M} = 0$.
  \begin{equation}
    P_{\text{eq}} = 2 P_L + 0 = 2 \left(\frac{1}{40}\right) = \frac{1}{20}\,\mathrm{cm}^{-1} \implies \mathbf{F_{\text{eq}} = -20\,\mathrm{cm}}
  \end{equation}
  Acts as a \textbf{concave mirror of focal length $20\,\mathrm{cm}$}.
  \item \textbf{Case (ii): Curved Surface is Silvered:}
  Concave mirror formed by silvered curved face has $R_M = 20\,\mathrm{cm} \implies f_M = -10\,\mathrm{cm} \implies P_M = -\frac{1}{f_M} = +\frac{1}{10}\,\mathrm{cm}^{-1}$.
  \begin{equation}
    P_{\text{eq}} = 2 P_L + P_M = 2 \left(\frac{1}{40}\right) + \frac{1}{10} = \frac{1}{20} + \frac{1}{10} = \frac{3}{20}\,\mathrm{cm}^{-1} \implies \mathbf{F_{\text{eq}} = -\frac{20}{3}\,\mathrm{cm} \approx -6.67\,\mathrm{cm}}
  \end{equation}
  Acts as a \textbf{concave mirror of focal length $6.67\,\mathrm{cm}$}.
\end{enumerate}
\end{usecasebox}

\begin{usecasebox}[Archetype 2: Multi-Wavelength Bright Fringe Coincidence in YDSE]
\textbf{Problem:} In a YDSE setup with slit separation $d = 1.0\,\mathrm{mm}$ and screen distance $D = 1.0\,\mathrm{m}$, light containing two wavelengths $\lambda_1 = 600\,\mathrm{nm}$ and $\lambda_2 = 480\,\mathrm{nm}$ is used. Determine the minimum non-zero distance from the central maximum where bright fringes of both wavelengths coincide.
\textbf{Step-by-Step Solution:}
\begin{enumerate}
  \item For bright fringes to coincide at position $y$:
  \begin{equation}
    y = \frac{n_1 \lambda_1 D}{d} = \frac{n_2 \lambda_2 D}{d} \implies n_1 \lambda_1 = n_2 \lambda_2
  \end{equation}
  \item Taking the ratio of integers:
  \begin{equation}
    \frac{n_1}{n_2} = \frac{\lambda_2}{\lambda_1} = \frac{480\,\mathrm{nm}}{600\,\mathrm{nm}} = \frac{4}{5}
  \end{equation}
  \item Smallest non-zero integers satisfying this condition: $n_1 = 4$ and $n_2 = 5$.
  \item Minimum distance $y_{\min}$:
  \begin{equation}
    y_{\min} = \frac{4 \times (600 \times 10^{-9}\,\mathrm{m}) \times 1.0\,\mathrm{m}}{1.0 \times 10^{-3}\,\mathrm{m}} = 2.4 \times 10^{-3}\,\mathrm{m} = \mathbf{2.4\,\mathrm{mm}}
  \end{equation}
  The 4th bright fringe of $\lambda_1$ coincides with the 5th bright fringe of $\lambda_2$ at $y = 2.4\,\mathrm{mm}$.
\end{enumerate}
\end{usecasebox}

\begin{usecasebox}[Archetype 3: Zener Diode Voltage Regulator Design Calculations]
\textbf{Problem:} A Zener diode with breakdown voltage $V_z = 10\,\mathrm{V}$ and maximum power dissipation $P_{z,\max} = 1.0\,\mathrm{W}$ is used to regulate a DC supply that fluctuates between $15\,\mathrm{V}$ and $25\,\mathrm{V}$. The load resistance is $R_L = 1.0\,\mathrm{k}\Omega$. Calculate: (i) the load current $I_L$, and (ii) the suitable series resistance $R_s$ to ensure the Zener current remains within safe limits.
\textbf{Step-by-Step Solution:}
\begin{enumerate}
  \item The regulated voltage across the load is fixed at $V_L = V_z = 10\,\mathrm{V}$.
  \item Constant load current:
  \begin{equation}
    I_L = \frac{V_z}{R_L} = \frac{10\,\mathrm{V}}{1000\,\Omega} = 10\,\mathrm{mA} = 0.01\,\mathrm{A}
  \end{equation}
  \item Maximum safe Zener diode current:
  \begin{equation}
    I_{z,\max} = \frac{P_{z,\max}}{V_z} = \frac{1.0\,\mathrm{W}}{10\,\mathrm{V}} = 100\,\mathrm{mA} = 0.10\,\mathrm{A}
  \end{equation}
  \item Total maximum supply current: $I_{\text{total},\max} = I_{z,\max} + I_L = 100\,\mathrm{mA} + 10\,\mathrm{mA} = 110\,\mathrm{mA}$.
  \item At maximum input voltage $V_{\text{in}} = 25\,\mathrm{V}$:
  \begin{equation}
    R_{s,\min} = \frac{V_{\text{in},\max} - V_z}{I_{\text{total},\max}} = \frac{25 - 10}{0.11} = \frac{15}{0.11} \approx \mathbf{136.4\,\Omega}
  \end{equation}
  \item At minimum input voltage $V_{\text{in}} = 15\,\mathrm{V}$ with minimum holding current $I_{z,\min} \approx 5\,\mathrm{mA}$:
  \begin{equation}
    R_{s,\max} = \frac{V_{\text{in},\min} - V_z}{I_{z,\min} + I_L} = \frac{15 - 10}{0.005 + 0.010} = \frac{5}{0.015} \approx \mathbf{333.3\,\Omega}
  \end{equation}
  A standard resistance of $R_s \approx 200\,\Omega$ provides robust regulation across the entire input range.
\end{enumerate}
\end{usecasebox}

\end{document}
